% GENERATED FILE. Edit canonical lessons, then run npm run generate:books.
% canonical-source-hash: fnv1a64:ba03905b69fad14e
% canonical-lessons: HI-C437-mansun, HI-C437-barsat, HI-C437-nami, HI-C437-chana, HI-C437-pasijna, HI-R437-review-430-437

\chapter{The Monsoon}
\label{ch:hi-the-monsoon}
\hlchaptermodality{437}

\begin{chapteropening}
\textbf{By the end of this chapter:} \emph{I can talk about the monsoon and the rainy season, damp walls, clouds spreading over and sweaty hands.}

Complete the last lesson of chapter 437: I can talk about the monsoon and the rainy season, damp walls, clouds spreading over and sweaty hands.
\end{chapteropening}

\section[\texorpdfstring{\hi{मानसून} (mānsūn)}{मानसून (mānsūn)}]{\hi{मानसून} (mānsūn) — the monsoon}

\label{lesson:HI-C437-mansun}



\begin{quote}
Before the new one: say the Hindi for a raincoat, then the Hindi for to pour (of rain).
\end{quote}

\subsection*{You'll want to know: \hi{मानसून}}
\textbf{\hi{मानसून}} — \emph{mānsūn} — \textquotedblleft{}the monsoon\textquotedblright{}.

\begin{quote}
\hi{इस} \hi{साल} \hi{मानसून} \hi{जल्दी} \hi{आया।}
\end{quote}

\begin{quote}
\emph{is sāl mānsūn jaldī āyā.}
\end{quote}

\begin{quote}
The monsoon came early this year.
\end{quote}

Masculine.

\begin{cousinweb}
\textbf{\hi{मानसून}} (\emph{mānsūn}) is borrowed from English \emph{monsoon}, a word English itself took from the Indian Ocean trade; Hindi also has its own word for the rainy season.
\end{cousinweb}

\begin{grammarlens}[title={What you've built}]
One more word for getting around.
\end{grammarlens}

\subsection*{Your turn}
\begin{itemize}
\raggedright
  \item \emph{Say it:} \emph{mānsūn}
  \item \emph{Say it:} \emph{mānsūn}, once more
  \item \emph{Say it:} \emph{is sāl mānsūn jaldī āyā.} — The monsoon came early this year.
  \item \emph{Recall:} say the Hindi for a raincoat, then the Hindi for to pour (of rain), then say \emph{mānsūn} again
\end{itemize}

\begin{tcolorbox}[breakable,colback=teal!4,colframe=teal!35!black,title={Before you move on}]
What does \hi{मानसून} mean? (\textquotedblleft{}The monsoon\textquotedblright{}.) Say it once more.
\end{tcolorbox}
\section[\texorpdfstring{\hi{बरसात} (barsāt)}{बरसात (barsāt)}]{\hi{बरसात} (barsāt) — the rainy season}

\label{lesson:HI-C437-barsat}



\begin{quote}
Before the new one: say the Hindi for to pour (of rain), then the Hindi for the monsoon.
\end{quote}

\subsection*{You'll want to know: \hi{बरसात}}
\textbf{\hi{बरसात}} — \emph{barsāt} — \textquotedblleft{}the rainy season\textquotedblright{}.

\begin{quote}
\hi{बरसात} \hi{में} \hi{गाँव} \hi{बहुत} \hi{हरा} \hi{हो} \hi{जाता} \hi{है।}
\end{quote}

\begin{quote}
\emph{barsāt meṁ gā̃v bahut harā ho jātā hai.}
\end{quote}

\begin{quote}
In the rainy season the village turns very green.
\end{quote}

Feminine.

\begin{cousinweb}
\textbf{\hi{बरसात}} (\emph{barsāt}) is a tadbhava of Sanskrit \emph{varṣārātri}, \textquotedblleft{}the rainy season\textquotedblright{} (CDIAL 11398), literally the rain nights; CDIAL lists Hindi \emph{barsāt}.
\end{cousinweb}

\begin{grammarlens}[title={What you've built}]
One more word for getting around.
\end{grammarlens}

\subsection*{Your turn}
\begin{itemize}
\raggedright
  \item \emph{Say it:} \emph{barsāt}
  \item \emph{Say it:} \emph{barsāt}, once more
  \item \emph{Say it:} \emph{barsāt meṁ gā̃v bahut harā ho jātā hai.} — In the rainy season the village turns very green.
  \item \emph{Recall:} say the Hindi for to pour (of rain), then the Hindi for the monsoon, then say \emph{barsāt} again
\end{itemize}

\begin{tcolorbox}[breakable,colback=teal!4,colframe=teal!35!black,title={Before you move on}]
What does \hi{बरसात} mean? (\textquotedblleft{}The rainy season\textquotedblright{}.) Say it once more.
\end{tcolorbox}
\section[\texorpdfstring{\hi{नमी} (namī)}{नमी (namī)}]{\hi{नमी} (namī) — damp; moisture}

\label{lesson:HI-C437-nami}



\begin{quote}
Before the new one: say the Hindi for the monsoon, then the Hindi for the rainy season.
\end{quote}

\subsection*{You'll want to know: \hi{नमी}}
\textbf{\hi{नमी}} — \emph{namī} — \textquotedblleft{}damp; moisture\textquotedblright{}.

\begin{quote}
\hi{बरसात} \hi{में} \hi{दीवारों} \hi{पर} \hi{नमी} \hi{आ} \hi{जाती} \hi{है।}
\end{quote}

\begin{quote}
\emph{barsāt meṁ dīvāroṁ par namī ā jātī hai.}
\end{quote}

\begin{quote}
In the rainy season the walls get damp.
\end{quote}

Feminine.

\begin{cousinweb}
\textbf{\hi{नमी}} (\emph{namī}) is built on \emph{nam}, \textquotedblleft{}damp, moist\textquotedblright{}, usually traced to Persian \emph{nam}, with the ending -ī that makes a noun.
\end{cousinweb}

\begin{grammarlens}[title={What you've built}]
One more word for getting around.
\end{grammarlens}

\subsection*{Your turn}
\begin{itemize}
\raggedright
  \item \emph{Say it:} \emph{namī}
  \item \emph{Say it:} \emph{namī}, once more
  \item \emph{Say it:} \emph{barsāt meṁ dīvāroṁ par namī ā jātī hai.} — In the rainy season the walls get damp.
  \item \emph{Recall:} say the Hindi for the monsoon, then the Hindi for the rainy season, then say \emph{namī} again
\end{itemize}

\begin{tcolorbox}[breakable,colback=teal!4,colframe=teal!35!black,title={Before you move on}]
What does \hi{नमी} mean? (\textquotedblleft{}Damp; moisture\textquotedblright{}.) Say it once more.
\end{tcolorbox}
\section[\texorpdfstring{\hi{छाना} (chānā)}{छाना (chānā)}]{\hi{छाना} (chānā) — to spread over (of clouds); to cover}

\label{lesson:HI-C437-chana}



\begin{quote}
Before the new one: say the Hindi for the rainy season, then the Hindi for damp.
\end{quote}

\subsection*{You'll want to know: \hi{छाना}}
\textbf{\hi{छाना}} — \emph{chānā} — \textquotedblleft{}to spread over (of clouds); to cover\textquotedblright{}.

\begin{quote}
\hi{आसमान} \hi{में} \hi{काले} \hi{बादल} \hi{छा} \hi{गए।}
\end{quote}

\begin{quote}
\emph{āsmān meṁ kāle bādal chā gae.}
\end{quote}

\begin{quote}
Dark clouds spread across the sky.
\end{quote}

Past (it …): \hi{वह} \hi{कोहरा} \hi{पूरे} \hi{शहर} \hi{पर} \hi{छा} \hi{गया}, \emph{vah kohrā pūre śahar par chā gayā}.

\begin{cousinweb}
\textbf{\hi{छाना}} (\emph{chānā}) is a tadbhava of Sanskrit \emph{chādayati}, \textquotedblleft{}covers, veils\textquotedblright{} (CDIAL 5018), which lists Hindi \emph{chānā}, \textquotedblleft{}to thatch, spread around\textquotedblright{}: clouds that \hi{छा} \hi{जाते} cover the sky like a roof.
\end{cousinweb}

\begin{grammarlens}[title={What you've built}]
One more word for getting around.
\end{grammarlens}

\subsection*{Your turn}
\begin{itemize}
\raggedright
  \item \emph{Say it:} \emph{chānā}
  \item \emph{Say it:} \emph{chānā}, once more
  \item \emph{Say it:} \emph{āsmān meṁ kāle bādal chā gae.} — Dark clouds spread across the sky.
  \item \emph{Recall:} say the Hindi for the rainy season, then the Hindi for damp, then say \emph{chānā} again
\end{itemize}

\begin{tcolorbox}[breakable,colback=teal!4,colframe=teal!35!black,title={Before you move on}]
What does \hi{छाना} mean? (\textquotedblleft{}To spread over (of clouds); to cover\textquotedblright{}.) Say it once more.
\end{tcolorbox}
\section[\texorpdfstring{\hi{पसीजना} (pasījnā)}{पसीजना (pasījnā)}]{\hi{पसीजना} (pasījnā) — to sweat; to soften (of the heart)}

\label{lesson:HI-C437-pasijna}



\begin{quote}
Before the new one: say the Hindi for damp, then the Hindi for to spread over.
\end{quote}

\subsection*{You'll want to know: \hi{पसीजना}}
\textbf{\hi{पसीजना}} — \emph{pasījnā} — \textquotedblleft{}to sweat; to soften (of the heart)\textquotedblright{}.

\begin{quote}
\hi{बरसात} \hi{की} \hi{गर्मी} \hi{में} \hi{हाथ} \hi{पसीजते} \hi{हैं।}
\end{quote}

\begin{quote}
\emph{barsāt kī garmī meṁ hāth pasījte haiṁ.}
\end{quote}

\begin{quote}
In the damp heat of the rains your hands get sweaty.
\end{quote}

Past (it …): \hi{वह} \hi{गिलास} \hi{ठंड} \hi{से} \hi{पसीज} \hi{गया}, \emph{vah gilās ṭhaṇḍ se pasīj gayā}.

\begin{cousinweb}
\textbf{\hi{पसीजना}} (\emph{pasījnā}) goes back to \emph{*prasvidyati}, \textquotedblleft{}sweats\textquotedblright{}, a form CDIAL reconstructs (CDIAL 8896), which lists Hindi \emph{pasījnā}; \hi{पसीना}, sweat, is its cousin. A heart that \hi{पसीजता} softens with pity.
\end{cousinweb}

\begin{grammarlens}[title={What you've built}]
One more word for getting around.
\end{grammarlens}

\subsection*{Your turn}
\begin{itemize}
\raggedright
  \item \emph{Say it:} \emph{pasījnā}
  \item \emph{Say it:} \emph{pasījnā}, once more
  \item \emph{Say it:} \emph{barsāt kī garmī meṁ hāth pasījte haiṁ.} — In the damp heat of the rains your hands get sweaty.
  \item \emph{Recall:} say the Hindi for damp, then the Hindi for to spread over, then say \emph{pasījnā} again
\end{itemize}

\begin{tcolorbox}[breakable,colback=teal!4,colframe=teal!35!black,title={Before you move on}]
What does \hi{पसीजना} mean? (\textquotedblleft{}To sweat; to soften (of the heart)\textquotedblright{}.) Say it once more.
\end{tcolorbox}
\section[(dialogue)]{Looking back: from Capital and Bankruptcy to The Monsoon}

\label{lesson:HI-R437-review-430-437}



\begin{quote}
Say the Hindi for bankruptcy and a fund.
\end{quote}

\subsection*{Your turn}
Say these aloud:

\begin{itemize}
\raggedright
  \item bankruptcy, a fund, capital (money), to be paid off, to collect (dues)
  \item the going rate, worth, a deal, a sale, good value
  \item goods, a store, a packet, a label, to wrap
  \item a queue, a token, a counter, a weekly market, to hand (to)
  \item real, pure, to rub, to suit, a touchstone
  \item loose coins, a size, an outfit, a handkerchief, to huddle up
  \item a drop, a cloud (literary), steam, a raincoat, to pour (of rain)
  \item the monsoon, the rainy season, damp, to spread over, to sweat
\end{itemize}

\begin{tcolorbox}[breakable,colback=teal!4,colframe=teal!35!black,title={Before you move on}]
Bankruptcy? (\textbf{\hi{दिवाला}}, \emph{divālā}.) A fund? (\textbf{\hi{निधि}}, \emph{nidhi}.) Capital (money)? (\textbf{\hi{पूँजी}}, \emph{pū̃jī}.) To be paid off? (\textbf{\hi{चुकना}}, \emph{cuknā}.) To collect (dues)? (\textbf{\hi{उगाहना}}, \emph{ugāhnā}.) The going rate? (\textbf{\hi{भाव}}, \emph{bhāv}.) Worth? (\textbf{\hi{मोल}}, \emph{mol}.) A deal? (\textbf{\hi{सौदा}}, \emph{saudā}.) A sale? (\textbf{\hi{सेल}}, \emph{sel}.) Good value? (\textbf{\hi{किफ़ायती}}, \emph{kifāyatī}.) Goods? (\textbf{\hi{माल}}, \emph{māl}.) A store? (\textbf{\hi{भंडार}}, \emph{bhaṇḍār}.) A packet? (\textbf{\hi{पैकेट}}, \emph{paikeṭ}.) A label? (\textbf{\hi{लेबल}}, \emph{lebal}.) To wrap? (\textbf{\hi{लपेटना}}, \emph{lapeṭnā}.) A queue? (\textbf{\hi{क़तार}}, \emph{qatār}.) A token? (\textbf{\hi{टोकन}}, \emph{ṭokan}.) A counter? (\textbf{\hi{काउंटर}}, \emph{kāunṭar}.) A weekly market? (\textbf{\hi{हाट}}, \emph{hāṭ}.) To hand (to)? (\textbf{\hi{थमाना}}, \emph{thamānā}.) Real? (\textbf{\hi{असली}}, \emph{aslī}.) Pure? (\textbf{\hi{शुद्ध}}, \emph{śuddh}.) To rub? (\textbf{\hi{रगड़ना}}, \emph{ragaṛnā}.) To suit? (\textbf{\hi{फबना}}, \emph{phabnā}.) A touchstone? (\textbf{\hi{कसौटी}}, \emph{kasauṭī}.) Loose coins? (\textbf{\hi{रेज़गारी}}, \emph{rezgārī}.) A size? (\textbf{\hi{साइज़}}, \emph{sāiz}.) An outfit? (\textbf{\hi{पोशाक}}, \emph{pośāk}.) A handkerchief? (\textbf{\hi{रुमाल}}, \emph{rumāl}.) To huddle up? (\textbf{\hi{सिमटना}}, \emph{simaṭnā}.) A drop? (\textbf{\hi{बूँद}}, \emph{bū̃d}.) A cloud (literary)? (\textbf{\hi{मेघ}}, \emph{megh}.) Steam? (\textbf{\hi{भाप}}, \emph{bhāp}.) A raincoat? (\textbf{\hi{रेनकोट}}, \emph{renkoṭ}.) To pour (of rain)? (\textbf{\hi{बरसना}}, \emph{barasnā}.) The monsoon? (\textbf{\hi{मानसून}}, \emph{mānsūn}.) The rainy season? (\textbf{\hi{बरसात}}, \emph{barsāt}.) Damp? (\textbf{\hi{नमी}}, \emph{namī}.) To spread over? (\textbf{\hi{छाना}}, \emph{chānā}.) To sweat? (\textbf{\hi{पसीजना}}, \emph{pasījnā}.)
\end{tcolorbox}
